%% kapa: exam=1 round=37 year=2026 period=2 form=A subject=accounting subject_ko=회계학 seq=25 number=65
%% source: Q-Net 감정평가사 기출문제 (cst00309) artlSeq=5258924 file=raw/1008/5258924/2026년 제37회 감정평가사 1차 2교시.pdf
%% conversion: kordoc 4.24.51 -> tools/gen1.py (auto); review=proofread-sonnet; flags=none
\begin{question}{65}
(주)감평은 20×1년 초 사용목적 기계장치(취득원가 ￦1,000,000, 잔존가치 ￦0, 내용연수 4년, 연수합계법 상각)를 취득하여 원가모형을 적용하고 있다. 20×2년 초 상기 자산에 대한 감가상각방법을 정액법으로 변경하였으며, 동 일자에 자본적 지출에 해당하는 ￦300,000을 지출(기계장치 장부금액 증가)하여 내용연수가 2년 연장되고 잔존가치는 ￦200,000으로 증가하였다. 동 기계장치 관련 회계처리가 (주)감평의 20×2년도 당기순이익에 미치는 영향은? (단, 단수차이로 인해 오차가 있으면 가장 근사치를 선택한다.)
\begin{choices}
  \item ￦116,667 감소
  \item ￦140,000 감소
  \item ￦183,333 감소
  \item ￦193,333 감소
  \item ￦220,000 감소
\end{choices}
\end{question}
\answer{2}
